% CLAIM: what was measured, on what, and how far it generalizes. Protocol constants live
% in the committed sweep configs; every result JSON stamps commit and environment.
\section{Setup}
\label{sec:setup}

\paragraph{Platforms.} We use one node of eight A100-SXM4-80GB GPUs connected by NVLink,
and one TPU v5e-8 host with eight 16\,GB chips connected by ICI, the chips' own
interconnect \citep{nvidiaA100,tpuv5e}. An A100 has 80\,GB of memory against 16\,GB for a
v5e chip, and the v5e's matrix unit multiplies in bfloat16 by default. The host-boundary
measurements use two A100 nodes connected by TCP sockets (Section~\ref{subsec:boundary}).

\paragraph{The block benchmark.} Most measurements use one simplified transformer block:
six square $d\times d$ matrices (attention q, k, v and o, and a two-layer MLP of width
$d$), one attention head, no normalization parameters, and a squared-error loss on
synthetic inputs. So $P = 6d^2$; at $d{=}2048$ that is 25.2M parameters and a 96\,MiB
float32 update. Each generation evaluates the whole population on one batch of eight
sequences of length 32. The placement grid covers $d\in\{512,2048\}$, two populations per
width (256 and 1024 at $d{=}512$; 128 and 256 at $d{=}2048$) and $D\in\{1,2,4,8\}$; the
cost grid of Appendix~\ref{sec:systems-ablations} is wider. The simplifications change how
much computation the block does per byte of memory traffic, and how much memory it needs,
relative to a full model. Its absolute ratios therefore apply only to this workload, and
the real model below tests whether the faster placement stays the same.

\paragraph{The real model.} We time complete generations of Qwen2.5-0.5B-Instruct
\citep{qwen25} under both placements on the v5e (Figure~\ref{fig:f9}); no A100 grid was
run for the real model. As on the block, low-rank candidates are evaluated in one batch and seed candidates
one at a time. The parameters, the reconstruction and the all-reduce use float32, and the
model is evaluated with a bfloat16 copy of the parameters. The evaluation procedure differs
between perturbation types, but never between the two placements being compared. The
fitness is the mean next-token negative log-likelihood on eight fixed Countdown prompts,
padded to 128 tokens, with padding excluded from the loss.

\paragraph{Protocol.} Complete-generation benchmarks use three warm-up generations
followed by five timed repeats, or seven for the host-boundary experiment. We report the
median and wait for device execution to finish before stopping each timer. Separate
component-timing procedures are described in Appendix~\ref{sec:implementation}.
Matrix-multiplication precision is fixed per experiment and recorded. The
scaling and placement sweeps on the block, and the separate reconstruction timings used
for Eq.~\ref{eq:crossover}, use float32 storage and JAX's highest precision on both
platforms; the cost grid uses the default precision, and
Appendix~\ref{sec:systems-ablations} compares the two settings. Both platforms run the
same configurations of model size, population, perturbation type and datatype. A
configuration that does not fit is recorded as out of memory (OOM) rather than shrunk to
fit, so memory limits remain part of the comparison.

\paragraph{Scope.} The measurements cover one host per platform, where communication
between devices is fastest, and one pair of A100 hosts connected by TCP sockets; one
simplified block; and one real model, Qwen2.5-0.5B, on one platform. Results are stated
within those limits.

\paragraph{Provenance.} Experiment configurations are committed before runs. Result
records capture the code versions and execution environments, and measurement figures
and tables are generated from saved results (Appendix~\ref{sec:repro}).
